% problem_id: S001-lemma-find-affine-codimension-1
% source_id: S001 (Stacks Project)
% source_file: more-groupoids.tex
% statement_locator: more-groupoids.tex lines 2805--2815
% proof_locator: more-groupoids.tex lines 2817--2950
% env: lemma
% id_note: label 在本来源内唯一
% pairing: tight (verified)
% status: complete (statement + author proof verbatim + reason + locators)
%
\begin{lemma}
\label{lemma-find-affine-codimension-1}
Let $(U, R, s, t, c)$ be a groupoid scheme. Let $u \in U$. Assume
\begin{enumerate}
\item $s, t$ are finite morphisms,
\item $U$ is separated and locally Noetherian,
\item $\dim(\mathcal{O}_{U, u'}) \leq 1$ for every point $u'$
in the orbit of $u$.
\end{enumerate}
Then $u$ is contained in an $R$-invariant affine open of $U$.
\end{lemma}

\begin{proof}
The $R$-orbit of $u$ is finite. By conditions (2) and (3) it is contained
in an affine open $U'$ of $U$, see
Varieties, Proposition
\ref{varieties-proposition-finite-set-of-points-of-codim-1-in-affine}.
Then $t(s^{-1}(U \setminus U'))$ is an $R$-invariant
closed subset of $U$ which does not contain $u$. Thus
$U \setminus t(s^{-1}(U \setminus U'))$ is an $R$-invariant open
of $U'$ containing $u$.
Replacing $U$ by this open we may assume $U$ is quasi-affine.

\medskip\noindent
By Lemma \ref{lemma-find-affine-integral} we may replace $U$ by its reduction
and assume $U$ is reduced. This means $R$-invariant subschemes
$W' \subset W \subset U$ of
Lemma \ref{lemma-finite-flat-over-almost-dense-subscheme}
are equal $W' = W$. As $U = t(s^{-1}(\overline{W}))$ some point
$u'$ of the $R$-orbit of $u$ is contained in $\overline{W}$
and by Lemma \ref{lemma-find-affine-integral}
we may replace $U$ by $\overline{W}$ and $u$ by $u'$.
Hence we may assume there is
a dense open $R$-invariant subscheme $W \subset U$ such that
the morphisms $s_W, t_W$ of the restriction $(W, R_W, s_W, t_W, c_W)$ are
finite locally free.

\medskip\noindent
If $u \in W$ then we are done by
Groupoids, Lemma \ref{groupoids-lemma-find-invariant-affine}
(because $W$ is quasi-affine so any finite set of points
of $W$ is contained in an affine open, see
Properties, Lemma \ref{properties-lemma-ample-finite-set-in-affine}).
Thus we assume $u \not \in W$ and hence none of the points of the
orbit of $u$ is in $W$. Let $\xi \in U$
be a point with a nontrivial specialization to a point $u'$ in the orbit
of $u$. Since there are no specializations among the points in the
orbit of $u$ (Lemma \ref{lemma-no-specializations-map-to-same-point})
we see that $\xi$ is not in the orbit.
By assumption (3) we see that $\xi$ is a generic point of $U$
and hence $\xi \in W$.
As $U$ is Noetherian there are finitely many of these
points $\xi_1, \ldots, \xi_m \in W$. Because $s_W, t_W$ are flat the orbit
of each $\xi_j$ consists of generic points of irreducible components
of $W$ (and hence $U$).

\medskip\noindent
Let $j : U \to \Spec(A)$ be an immersion of $U$ into an affine scheme
(this is possible as $U$ is quasi-affine). Let $J \subset A$
be an ideal such that $V(J) \cap j(W) = \emptyset$ and $V(J) \cup j(W)$
is closed. Apply Lemma \ref{lemma-find-almost-invariant-function}
to the groupoid scheme $(W, R_W, s_W, t_W, c_W)$, the morphism
$j|_W : W \to \Spec(A)$, the points $\xi_j$, and the ideal $J$
to find an $f \in J$ such that $(j|_W)^{-1}D(f)$ is an $R_W$-invariant
affine open containing $\xi_j$ for all $j$. Since $f \in J$
we see that $j^{-1}D(f) \subset W$, i.e., $j^{-1}D(f)$ is
an $R$-invariant affine open of $U$ contained in $W$
containing all $\xi_j$.

\medskip\noindent
Let $Z$ be the reduced induced closed subscheme structure on
$$
U \setminus j^{-1}D(f) = j^{-1}V(f).
$$
Then $Z$ is set theoretically
$R$-invariant (but it may not be scheme theoretically $R$-invariant).
Let $(Z, R_Z, s_Z, t_Z, c_Z)$ be the restriction of $R$ to $Z$.
Since $Z \to U$ is finite, it follows that $s_Z$ and $t_Z$ are finite.
Since $u \in Z$ the orbit of $u$ is in $Z$ and agrees with the
$R_Z$-orbit of $u$ viewed as a point of $Z$. Since
$\dim(\mathcal{O}_{U, u'}) \leq 1$ and since $\xi_j \not \in Z$
for all $j$, we see that $\dim(\mathcal{O}_{Z, u'}) \leq 0$ for
all $u'$ in the orbit of $u$. In other words, the $R_Z$-orbit of $u$
consists of generic points of irreducible components of $Z$.

\medskip\noindent
Let $I \subset A$ be an ideal such that $V(I) \cap j(U) =\emptyset$
and $V(I) \cup j(U)$ is closed. Apply
Lemma \ref{lemma-find-almost-invariant-function} to
the groupoid scheme $(Z, R_Z, s_Z, t_Z, c_Z)$, the restriction $j|_Z$,
the ideal $I$, and the point $u \in Z$ to obtain $h \in I$ such that
$j^{-1}D(h) \cap Z$ is an $R_Z$-invariant open affine containing $u$.

\medskip\noindent
Consider the $R_W$-invariant (Groupoids, Lemma
\ref{groupoids-lemma-determinant-trick}) function
$$
g = 
\text{Norm}_{s_W}(t_W^\sharp(j^\sharp(h)|_W)) \in \Gamma(W, \mathcal{O}_W)
$$
(In the following we only need the restriction of $g$ to $j^{-1}D(f)$ and
in this case the norm is along a finite locally free morphism of affines.)
We claim that
$$
V = (W_g \cap j^{-1}D(f)) \cup (j^{-1}D(h) \cap Z)
$$
is an $R$-invariant affine open of $U$ which finishes the proof of the lemma.
It is set theoretically $R$-invariant by construction. As $V$ is a
constructible set, to see that it is open it suffices to show it is
closed under generalization in $U$ (Topology, Lemma
\ref{topology-lemma-characterize-closed-Noetherian}
or the more general
Topology, Lemma
\ref{topology-lemma-constructible-stable-specialization-closed}).
Since $W_g \cap j^{-1}D(f)$ is open in $U$, it suffices to consider
a specialization $u_1 \leadsto u_2$ of $U$ with
$u_2 \in j^{-1}D(h) \cap Z$.
This means that $h$ is nonzero in $j(u_2)$ and $u_2 \in Z$.
If $u_1 \in Z$, then $j(u_1) \leadsto j(u_2)$ and since
$h$ is nonzero in $j(u_2)$ it is nonzero in $j(u_1)$ which
implies $u_1 \in V$. If $u_1 \not \in Z$ and
also not in $W_g \cap j^{-1}D(f)$, then $u_1 \in W$, $u_1 \not \in W_g$
because the complement of $Z = j^{-1}V(f)$ is contained in $W \cap j^{-1}D(f)$.
Hence there exists a point $r_1 \in R$ with $s(r_1) = u_1$
such that $h$ is zero in $t(r_1)$. Since $s$ is finite we
can find a specialization $r_1 \leadsto r_2$ with $s(r_2) = u_2$.
However, then we conclude that $h$ is zero in $u'_2 = t(r_2)$
which contradicts the fact that $j^{-1}D(h) \cap Z$
is $R$-invariant and $u_2$ is in it. Thus $V$ is open.

\medskip\noindent
Observe that $V \subset j^{-1}D(h)$ for our function $h \in I$.
Thus we obtain an immersion
$$
j' : V \longrightarrow \Spec(A_h)
$$
Let $f' \in A_h$ be the image of $f$. Then $(j')^{-1}D(f')$
is the principal open determined by $g$ in the affine
open $j^{-1}D(f)$ of $U$.
Hence $(j')^{-1}D(f)$ is affine. Finally,
$j'(V) \cap V(f') = j'(j^{-1}D(h) \cap Z)$
is closed in $\Spec(A_h/(f')) = \Spec((A/f)_h) = D(h) \cap V(f)$
by our choice of $h \in I$ and the ideal $I$. Hence we can apply
Lemma \ref{lemma-get-affine}
to conclude that $V$ is affine as claimed above.
\end{proof}
